% ============================================================================
% Section VIII: Conclusions and Future Work
% ============================================================================

\section{Conclusions}
\label{sec:conclusions}

This paper has presented GTOpt, an open-source optimizer for planning
and operating energy systems.  The main conclusions are:

\begin{enumerate}
  \item \textbf{One model across time scales.}  A single formulation of
    electricity, water, fuels and LNG, hydrogen, ammonia and heat is
    used from 15-minute operation to 30-year planning.  The scales are
    coupled through SDDP optimality cuts (water values and storage
    opportunity costs) and feasibility cuts (reservoir and water-right
    limits), so that the results of the different horizons are
    consistent by construction.

  \item \textbf{A broad formulation.}  The model covers the network
    with Kirchhoff constraints, switching and line losses, storage and
    hydro cascades with water rights, unit commitment, reserves and
    inertia, multi-carrier chains, emissions with caps and banking, and
    expansion of generation, transmission, storage and demand.  Operating
    rules beyond the element library are added through PAMPL.

  \item \textbf{Methods for large systems.}  The monolithic LP/MIP, SDDP
    with apertures, multi-cut and feasibility cuts, the multi-level
    cascade and the investment master, together with parallel
    subproblems, compressed snapshots and solver-neutral warm starts,
    make national-scale studies tractable on a single server.

  \item \textbf{Validated and auditable.}  GTOpt reproduces IEEE
    benchmarks and pandapower DC OPF to solver tolerance and produces
    valid bounds on published SDDP benchmarks.  On the Chilean daily
    schedule, the comparison with the CEN's PLEXOS solutions shows that
    the line-loss model drives the largest differences, and that a
    tangent loss model with a small loss price removes fictitious
    losses.  Exportable LP files, structural fingerprints and a
    published formulation make each result traceable.
\end{enumerate}

\subsection{Future Work}

GTOpt is the core of a larger set of modules being developed with the
user community under the GTOpt Community organization:

\begin{itemize}
  \item \textbf{Bid-based markets} (\texttt{gtopt-bidmark}): clearing
    offers instead of audited costs, in line with the transition of the
    Chilean power system.  The goal is one dispatch chain from
    medium-term water values to day-ahead and intraday scheduling, with
    co-optimized energy and ancillary services, nodal prices, capacity
    and long-term energy products, financial transmission rights, and a
    market simulator to test designs and bidding behaviour.
  \item \textbf{An observatory of the Chilean market}
    (\texttt{gtopt-obs-cl}): tools to observe, reproduce and analyse
    the daily dispatch of the national system, building on the
    comparison of Section~\ref{sec:cases}.
  \item \textbf{More reference cases} (\texttt{gtopt-cases}): the
    repository already covers six countries
    (Section~\ref{sec:gtopt_cases}); it will grow with expansion cases
    and with the power systems of other regions contributed by the
    community, as open reference data for validation and studies.
  \item \textbf{Closing the remaining gaps with the official schedule}:
    a systematic attribution of the marginal-cost differences of
    Section~\ref{sec:cases} to network, conversion and degeneracy
    effects, and recommendations on loss modeling under surplus.
  \item \textbf{Exact-expectation SDDP}: a backward pass that evaluates
    the exact expectation over the uncertainty tree, alongside the
    aperture approximation, and stronger cuts for integer decisions.
  \item \textbf{AC power flow}: second-order cone and linear
    approximations of reactive power and voltage for operation studies.
  \item \textbf{GPU and machine learning}: wider use of GPU solvers,
    and learned warm starts, scenario reduction and cut selection for
    SDDP.
\end{itemize}

GTOpt is intended to support reproducible research and transparent
studies of power systems.  Its source code, formulation and test cases
are public, and contributions of networks, case studies and code are
welcome.
